\section{硬件缺口补全：从结构到 OS 约束}

前面已经给出从门电路到 CPU 的路线，但还不够。只讲“CPU、cache、DMA、中断、设备”这些名词，读者仍然不知道它们为什么必须这样接，为什么 OS 要保存某些状态，为什么驱动要按奇怪的顺序写寄存器，为什么同一条 \asmcode{mov rax, qword ptr [rbx + 16]} 有时只是读内存，有时会变成缺页、MMIO 访问或设备错误。本节把这些缺口补上。

读每个硬件对象都用同一个模板：

\[
\text{物理/逻辑结构}
\rightarrow
\text{内部状态}
\rightarrow
\text{时序边界}
\rightarrow
\text{对上接口}
\rightarrow
\text{OS 不变量}
\rightarrow
\text{失败证据}
\]

如果一个解释没有说清内部状态和 OS 不变量，就还没有讲明白。例如“中断用于通知 CPU”不够，必须继续问：pending 位在哪里，mask 位怎样生效，in-service 位何时清，EOI 早了会怎样，设备状态没清会怎样，vector 如何绑定 CPU。“DMA 用于设备搬运内存”也不够，必须继续问：descriptor 谁写，doorbell 何时写，设备用什么地址，IOMMU 映射何时撤销，completion 前 buffer 能不能释放。

\subsection{总图：硬件不是平面盒子，而是多条状态链}

下面这张图把 OS 最常踩的硬件链路放在同一页。左侧是 CPU 内部提交链，右侧是内存翻译和 cache 链，下侧是设备 I/O 链。每条链都有自己的状态机；OS 做的事就是把这些状态机收敛成进程、页、文件、socket、request、completion 等软件对象。

\begin{pdfdiagram}[x=1cm,y=1cm]
  \node[os control, minimum width=1.55cm] (isa) at (0,1.25) {ISA};
  \node[os compute, minimum width=1.55cm] (front) at (1.9,1.25) {fetch/decode};
  \node[os compute, minimum width=1.55cm] (exec) at (3.8,1.25) {execute};
  \node[os state, minimum width=1.55cm] (retire) at (5.7,1.25) {retire};
  \node[os control, minimum width=1.55cm] (trap) at (7.6,1.25) {trap/IRQ};
  \node[os control, minimum width=1.55cm] (kernel) at (9.5,1.25) {kernel};

  \node[os bit, minimum width=1.25cm] (va) at (0,-0.4) {VA};
  \node[os memory, minimum width=1.55cm] (tlb) at (1.9,-0.4) {TLB};
  \node[os memory, minimum width=1.55cm] (pt) at (3.8,-0.4) {page walk};
  \node[os bit, minimum width=1.25cm] (pa) at (5.7,-0.4) {PA};
  \node[os memory, minimum width=1.55cm] (cache) at (7.6,-0.4) {cache};
  \node[os memory, minimum width=1.55cm] (dram) at (9.5,-0.4) {DRAM};

  \node[os control, minimum width=1.55cm] (drv) at (0,-2.05) {driver};
  \node[os memory, minimum width=1.55cm] (ring) at (1.9,-2.05) {ring};
  \node[os control, minimum width=1.55cm] (db) at (3.8,-2.05) {doorbell};
  \node[os block, minimum width=1.55cm] (dev) at (5.7,-2.05) {device};
  \node[os memory, minimum width=1.55cm] (iommu) at (7.6,-2.05) {IOMMU};
  \node[os state, minimum width=1.55cm] (comp) at (9.5,-2.05) {completion};

  \draw[os strong bus] (isa) -- (front);
  \draw[os strong bus] (front) -- (exec);
  \draw[os strong bus] (exec) -- (retire);
  \draw[os feedback] (retire) -- (trap);
  \draw[os strong bus] (trap) -- (kernel);

  \draw[os strong bus] (va) -- (tlb);
  \draw[os bus] (tlb) -- node[os label, above] {miss} (pt);
  \draw[os strong bus] (pt) -- (pa);
  \draw[os strong bus] (pa) -- (cache);
  \draw[os strong bus] (cache) -- (dram);
  \draw[os feedback] (pt) |- node[os label, right] {\#PF} (kernel);

  \draw[os strong bus] (drv) -- (ring);
  \draw[os strong bus] (ring) -- (db);
  \draw[os strong bus] (db) -- (dev);
  \draw[os strong bus] (dev) -- (iommu);
  \draw[os strong bus] (iommu) -- (comp);
  \draw[os feedback] (comp) |- node[os label, right] {MSI-X} (kernel);
\end{pdfdiagram}
\diagramnote{三条关键链：CPU 提交链、地址翻译与 cache 链、设备 DMA 链。OS 必须把它们的状态边界对齐。}

\subsection{时钟、复位和时钟域：状态为什么不能随便读}

硬件状态由时钟推进。一个模块没有时钟，寄存器不会采样；一个模块还在复位，内部寄存器可能保持默认值；两个模块处于不同 clock domain，同一事件穿越边界时需要同步器或握手协议。OS 驱动里的“先开电源、再开时钟、再解除复位、再写寄存器”来自这些物理约束。

\begin{lstlisting}
device_resume_order:
  enable parent power domain
  enable regulator
  enable functional clock
  wait clock stable or PLL locked
  deassert reset
  wait device ready bit
  program registers
  enable DMA and interrupts
\end{lstlisting}

如果顺序错了，错误表现常常不是立即崩溃，而是随机超时、寄存器读回全 1、DMA 不完成、中断不来。原因是 OS 以为自己在访问“寄存器变量”，实际访问的是一个还没进入 READY 状态的硬件状态机。

\[
state_{device}=READY
\Leftarrow
power\_good \land clock\_stable \land reset=0 \land init\_done
\]

\begin{longtable}{p{0.20\linewidth}p{0.34\linewidth}p{0.34\linewidth}}
\toprule
对象 & 硬件事实 & OS 约束 \\
\midrule
PLL/clock & 频率稳定需要 lock 时间 & 切频或启用后等待稳定证据 \\
reset line & 复位会清内部队列和寄存器 & reset 后重建 BAR 寄存器、队列和中断配置 \\
clock gate & 省电时关闭模块时钟 & 访问 MMIO 前必须 runtime resume \\
power domain & 断电后状态可能丢失 & suspend 前停止 DMA，resume 后重初始化 \\
watchdog & 独立时钟检测系统卡死 & 内核要定期喂狗，panic/死锁时保留证据 \\
\bottomrule
\end{longtable}

电源管理、热管理、设备 reset 不是外围章节。它们决定“寄存器是否有意义、DMA 是否还会发生、中断是否还能投递”。这也是为什么驱动不能只写 happy path。

\subsection{x86-64 指令从字节到提交：一条 add 的完整链路}

以 \asmcode{add rax, rbx} 为例。汇编文本先由汇编器编码成字节；CPU 用 RIP 从 I-cache 取字节；译码器识别操作、源寄存器和目的寄存器；寄存器文件读出旧 \code{rax/rbx}；ALU 做加法；结果先进入内部结果路径或重排序缓冲，最后在提交边界更新软件可见的 \code{rax} 和 RFLAGS。

\begin{lstlisting}
add rax, rbx:
  bytes = fetch(RIP)
  uop = decode(bytes)
  src0 = read_arch_or_renamed(rax)
  src1 = read_arch_or_renamed(rbx)
  tmp  = integer_adder(src0, src1)
  flags = compute_add_flags(src0, src1, tmp)
  retire in program order:
    rax = tmp
    RFLAGS = flags
    RIP = RIP + instruction_length
\end{lstlisting}

关键点是“执行”和“提交”不同。乱序 CPU 可以提前执行后面的微操作，但 OS 能依赖的是提交后的架构状态。异常入口保存的也是提交边界上的状态，而不是 CPU 内部临时物理寄存器。

\[
architectural\ state_{t+1}
= retire(architectural\ state_t,\ committed\ instruction)
\]

这条式子解释上下文切换的范围：OS 保存通用寄存器、RIP、RSP、RFLAGS、FPU/SIMD 等架构状态；不保存 ROB、调度队列、物理寄存器池、分支预测器每个内部表项。后者影响性能和侧信道，前者决定程序语义。

\subsection{一条 load 为什么会牵动 MMU、cache 和异常}

再看 \asmcode{mov rax, qword ptr [rbx + 16]}。它表面是一次内存读，内部至少有地址生成、canonical 检查、TLB 查找、PTE 权限检查、cache 查找、可能的 page walk、可能的 page fault、最后写回寄存器。

\begin{lstlisting}
mov rax, qword ptr [rbx + 16]:
  va = rbx + 16
  check canonical form
  if DTLB misses:
    pte = hardware_page_walk(CR3, va)
    if pte invalid or permission denied:
      raise #PF with CR2=va and error_code
    insert translation into DTLB
  pa = translate(va)
  line = data_cache_lookup(pa)
  if miss:
    fill line from lower cache or memory
  rax = select_qword(line, pa.offset)
\end{lstlisting}

把它写成时序链：

\[
\begin{aligned}
VA &= R_{base}+disp,\\
VPN &= \lfloor VA / P \rfloor,\quad offset=VA\bmod P,\\
PTE &= walk(CR3, VPN),\\
allow &= present \land permission(PTE, CPL, access),\\
PA &= PTE.PPN \times P + offset.
\end{aligned}
\]

若 \(allow=false\)，硬件不能把这条 load 当作完成。它要产生精确异常：

\[
\#PF = \{CR2=VA,\ RIP=RIP_{faulting},\ error\_code,\ CPL\}
\]

这就是虚拟内存的根。OS 不是在任意时刻“发现内存不够”，而是在硬件给出 faulting address、错误码和 faulting RIP 后，查 VMA、分配页、装 PTE、刷新必要 TLB，然后让同一条指令重试。没有精确异常，按需分页和 COW 都不成立。

\subsection{MMU 不只是翻译地址，也是在执行权限逻辑}

页表项不是单纯的“虚拟页到物理页映射”。它同时编码 present、writable、user/supervisor、execute disable、accessed、dirty、global、cache attribute 等信息。硬件 page walk 读取 PTE 后，做的是权限判定。

\[
\begin{aligned}
allow(access)={}&present \\
&{}\land (CPL=0 \lor user=1)\\
&{}\land (access\ne write \lor writable=1)\\
&{}\land (access\ne execute \lor NX=0)
\end{aligned}
\]

\begin{longtable}{p{0.18\linewidth}p{0.34\linewidth}p{0.36\linewidth}}
\toprule
PTE 位 & 硬件含义 & OS 用途 \\
\midrule
present & 翻译是否有效 & 懒分配、swap、非法地址判断 \\
writable & 写是否允许 & COW、\code{mprotect}、只读代码段 \\
user & Ring 3 是否可访问 & 用户/内核隔离 \\
NX & 是否禁止取指执行 & W\^{}X、栈不可执行、安全隔离 \\
accessed/dirty & 硬件记录访问/写入痕迹 & 页面回收、writeback、老化算法 \\
global/PCID & 控制 TLB 保留范围 & 降低地址空间切换成本 \\
cache attr & 控制内存类型 & RAM、MMIO、write-combining、PMem \\
\bottomrule
\end{longtable}

因此“特权”不是口号。Ring 3 不能访问内核页，是因为 PTE 的 user 位、CPL 检查和异常入口共同工作。内核不能把用户指针当可信地址，也是因为当前 CPL 和访问路径会决定 fault 语义。系统调用里 \code{copy\_from\_user} 的价值，就是把“可能 fault 的用户访问”放到可恢复边界内。

\subsection{TLB shootdown 的精确条件}

页表在内存里，TLB 在 CPU 里。修改 PTE 后，旧翻译可能仍在一个或多个 CPU 的 TLB 中。OS 不能只改页表就释放物理页。

\[
reuse(PA)
\Rightarrow
\forall cpu\in users(mm):\ old\_translation(cpu,VA\rightarrow PA)=false
\]

安全 unmap 顺序是：

\begin{lstlisting}
safe_unmap(mm, va):
  acquire page table lock
  clear or restrict PTE
  record affected CPUs for this mm
  release page table lock
  invalidate local TLB entry
  send shootdown IPI to remote CPUs in cpumask
  wait for acknowledgements
  now old physical page may be freed or reused
\end{lstlisting}

如果省略等待，旧 CPU 可能继续用旧 TLB 写这个物理页；而该页也许已经被分配给另一个进程、内核 slab 或 page cache。结果不是普通崩溃，而是越权写和随机数据破坏。这就是 shootdown 贵但不能省的原因。

\subsection{内存控制器、DRAM bank 和 NUMA：物理页也有位置}

DRAM 不是平面数组。地址会被内存控制器拆到 channel、rank、bank、row、column。一次访问可能经历 row activate、column access、precharge、refresh 和队列等待。多 socket 或多内存控制器系统还会形成 NUMA：同一个物理页离不同 CPU 的距离不同。

\[
T_{dram}=T_{queue}+T_{cmd}+T_{row}+T_{col}+T_{transfer}
\]

\[
BW_{max}\approx channels \times transfers/s \times bytes/transfer
\]

\begin{longtable}{p{0.18\linewidth}p{0.34\linewidth}p{0.36\linewidth}}
\toprule
硬件结构 & 带来的现象 & OS 设计后果 \\
\midrule
channel & 并行内存通路 & 内存带宽和插条位置相关 \\
bank/row buffer & 行命中比行冲突快 & 顺序访问、页着色和访问模式会影响延迟 \\
refresh & DRAM 周期性恢复电荷 & 极端延迟抖动和可靠性测试要考虑 \\
ECC & 检测/纠正位错误 & MCE、EDAC、页下线、DIMM 诊断 \\
NUMA node & CPU 到内存距离不同 & first-touch、迁移、亲和性、远端访问成本 \\
\bottomrule
\end{longtable}

调度器和页分配器为什么要知道 NUMA？因为把线程调到远端 CPU，或者把热页放在远端内存，会让每次 cache miss 多走更长互联路径。大机器上的性能不是单靠算法复杂度决定，也由“线程在哪个 CPU、页在哪个 node、设备 DMA 到哪个 node”决定。

\subsection{cache 一致性只覆盖普通内存，不自动覆盖所有设备语义}

多核 CPU 的一致性协议让普通 write-back RAM 的 cache line 在 CPU 间维持一致解释。但这不等于所有观察者都按程序顺序看到所有写入，也不等于设备寄存器可当普通变量访问。

\[
coherence: \text{same address order}
\qquad
ordering: \text{different address visibility order}
\]

三个常见误解必须拆开：

\begin{longtable}{p{0.26\linewidth}p{0.31\linewidth}p{0.31\linewidth}}
\toprule
误解 & 错在哪里 & 正确边界 \\
\midrule
cache 一致性等于屏障 & 一致性多约束同一 cache line & 发布对象、DMA descriptor、锁释放还要内存序 \\
store 返回等于设备看见 & store buffer 和互联可能延迟可见性 & doorbell 前用 DMA/MMIO 屏障表达顺序 \\
MMIO 是内存变量 & 设备寄存器有副作用和访问宽度要求 & 用 \code{readl/writel} 等 API 和页属性 \\
\bottomrule
\end{longtable}

设备发布请求的基本不变量是：

\[
descriptor\ fields\ visible
\prec_{device}
doorbell\ visible
\]

完成回收的基本不变量是：

\[
free(buffer)\Rightarrow completion\ observed \land dma\_unmap\ done \land owner=CPU
\]

这两个式子比“加一个屏障就好了”更有解释力。屏障只是在具体架构和总线上实现这些可见性关系的工具。

\subsection{APIC、vector 和中断入口：谁把事件交给哪个 CPU}

x86-64 上中断不是“设备调用内核函数”。设备或控制器产生事件，APIC 体系把事件映射成 vector，CPU 在可接收边界进入 IDT 指定入口。local APIC 管本 CPU 的 timer、IPI、EOI 和 vector 接收；IOAPIC 管传统外部 IRQ 重定向；MSI/MSI-X 让 PCIe 设备写一条消息触发 vector。

\begin{pdfdiagram}[x=1cm,y=1cm]
  \node[os block, minimum width=1.55cm] (dev) at (0,0.6) {device};
  \node[os block, minimum width=1.65cm] (msix) at (2.1,0.6) {MSI-X msg};
  \node[os control, minimum width=1.65cm] (apic) at (4.2,0.6) {local APIC};
  \node[os control, minimum width=1.45cm] (cpu) at (6.25,0.6) {CPU};
  \node[os state, minimum width=1.45cm] (idt) at (8.15,0.6) {IDT};
  \node[os control, minimum width=1.65cm] (entry) at (10.1,0.6) {handler};

  \node[os block, minimum width=1.55cm] (timer) at (0,-1.0) {timer};
  \node[os block, minimum width=1.55cm] (ipi) at (2.1,-1.0) {IPI};

  \draw[os strong bus] (dev) -- (msix);
  \draw[os strong bus] (msix) -- node[os label, above] {vector} (apic);
  \draw[os strong bus] (apic) -- (cpu);
  \draw[os strong bus] (cpu) -- (idt);
  \draw[os strong bus] (idt) -- (entry);
  \draw[os bus] (timer) -- (apic);
  \draw[os bus] (ipi) -- (apic);
  \draw[os feedback] (entry.south) .. controls +(0,-1.1) and +(0,-1.1) .. node[os label, below] {EOI} (apic.south);
\end{pdfdiagram}
\diagramnote{中断投递链。设备事件先变成 vector，再由 CPU 查 IDT 进入入口；handler 完成后还要满足设备确认和 APIC EOI。}

中断控制器至少维护三类状态：

\begin{lstlisting}
interrupt_state:
  pending[source]     // event has arrived
  masked[source]      // software blocks delivery
  in_service[vector]  // CPU accepted but not EOIed
\end{lstlisting}

顺序错误会产生具体故障：

\begin{longtable}{p{0.26\linewidth}p{0.30\linewidth}p{0.32\linewidth}}
\toprule
错误顺序 & 现象 & 原因 \\
\midrule
未清设备状态就 EOI & 中断风暴 & APIC 可再投递，设备条件仍为真 \\
清设备状态但忘 EOI & 后续中断不来 & vector 仍 in-service \\
过早开中断 & handler 重入破坏共享状态 & 内核不变量还没恢复 \\
所有队列同一 vector & 单 CPU 成为瓶颈 & 多队列硬件没有分散到多核 \\
\bottomrule
\end{longtable}

所以中断处理不是“快点执行函数”。它是设备状态机、APIC 状态机和内核并发状态共同收敛。

\subsection{定时器：抢占和超时来自可编程 deadline}

定时器硬件最小模型是计数器加比较器：

\[
interrupt = enable \land (counter \ge compare)
\]

周期 tick 是每次设定固定间隔；高精度 timer 是把 compare 设成最近 deadline；tickless idle 是在没有近期事件时把下一次中断推远。

\begin{lstlisting}
timer_reprogram(now):
  next = min(scheduler_deadline,
             hrtimer_deadline,
             rcu_deadline,
             device_timeout)
  local_apic_timer.compare = next
\end{lstlisting}

OS 里的超时不是后台线程“经常看看”。它依赖硬件在未来某个时刻把控制权交回内核。抢占也一样：若用户程序死循环，只有 timer interrupt 能强制进入内核，设置 need\_resched，并在返回路径切换任务。

\[
preempt(current)\Rightarrow timer\ IRQ \land kernel\ entry \land scheduler\ decision
\]

这解释了为什么禁中断区必须短。禁中断期间，timer 不能及时打断本 CPU，调度延迟、RCU 延迟、设备 completion 延迟都会被拉长。

\subsection{PCIe、BAR 和配置空间：设备不是凭空出现}

PCIe 设备先通过配置空间被发现。配置空间提供 vendor/device ID、class code、BAR、MSI-X capability、AER capability、DMA 能力等。BAR 不是寄存器本身，而是“设备申请一段 MMIO 地址窗口”的描述。OS 分配地址后，CPU 对该物理窗口的 load/store 会被 root complex 路由到设备。

\begin{lstlisting}
pci_probe_contract:
  read config space
  allocate and enable BAR resources
  set bus mastering only after DMA is prepared
  choose DMA mask and attach IOMMU domain
  allocate MSI-X vectors and set affinity
  initialize device registers and queues
  publish device to upper layer
\end{lstlisting}

关键不变量：

\[
bus\ mastering=enabled
\Rightarrow
IOMMU/domain\ ready \land DMA\ mask\ valid \land driver\ owns\ queues
\]

如果过早启用 bus mastering，设备可能在 OS 没有准备好映射和缓冲区时发起 DMA。可靠驱动在 probe 中不是“尽快开始工作”，而是先建立设备能访问什么、不能访问什么、完成后怎样通知的边界。

\subsection{NVMe 和网卡：同一种队列型设备骨架}

NVMe、网卡、virtio 和很多加速器都使用相似骨架：submission ring、completion ring、doorbell、MSI-X、DMA buffer、request tag。区别在命令格式和上层语义，不在基本状态机。

\begin{lstlisting}
queue_device_submit(req):
  tag = allocate_tag()
  map buffers for DMA
  fill descriptor[tag]
  store_release(descriptor.valid, true)
  dma_wmb()
  writel(new_tail, doorbell)

queue_device_complete(cqe):
  req = lookup_tag(cqe.tag)
  read status and bytes_done
  dma_unmap(req.buffers)
  complete upper object
  free tag
\end{lstlisting}

队列深度 \(Q\)、平均服务时间 \(S\) 和吞吐近似关系：

\[
throughput \le \min(device\_rate,\ Q/S)
\]

队列太浅，设备空等；队列太深，延迟和取消成本上升。中断合并可以减少 IRQ 次数，但会增加单个请求完成可见的延迟：

\[
latency_{observed}=latency_{device}+latency_{coalescing}+latency_{scheduler}
\]

这就是 blk-mq、NAPI、RSS、MSI-X affinity 都围绕队列和 CPU 亲和性设计的原因。设备硬件已经多队列，OS 必须把队列、CPU、NUMA node、软中断和上层对象匹配起来。

\subsection{USB 和低速总线：不是所有设备都像 PCIe}

PCIe 设备能较好自描述，USB 设备通过描述符和端点自描述；I2C、SPI、GPIO、PWM、ADC 这类板级硬件通常不能完整自描述，需要 ACPI 或板级描述告诉 OS 它们存在、地址是多少、中断接在哪个 GPIO、复位脚是哪根。

\begin{longtable}{p{0.18\linewidth}p{0.34\linewidth}p{0.36\linewidth}}
\toprule
总线/接口 & 硬件特点 & OS 约束 \\
\midrule
USB & 主机控制，端点和 URB 异步完成，热插拔频繁 & 断开后 fd 仍可能存在，操作要返回错误并释放资源 \\
I2C/SMBus & 低速共享总线，设备地址有限，事务短 & 访问要串行化，睡眠上下文限制，ACPI 描述很重要 \\
SPI & 片选线选择设备，控制器发起传输 & 片选、模式、速率、DMA/PIO 路径都要配置 \\
GPIO & 单 bit 输入输出，也可当中断线 & 方向、上下拉、边沿触发、去抖都影响语义 \\
PWM/ADC & 模拟世界和数字控制边界 & 精度、采样率、占空比、校准和权限都要处理 \\
\bottomrule
\end{longtable}

所以“驱动”不是一个统一写寄存器模型。USB 驱动管理 URB 生命周期；I2C 驱动管理低速寄存器事务；GPIO 驱动管理 pin 状态和中断边沿；PCIe 驱动管理 BAR、DMA 和 MSI-X。OS 统一它们的方式，是把不同硬件收敛成字符设备、输入事件、网络设备、块设备、hwmon、DRM object 等上层对象。

\subsection{显示、音频和实时流：为什么 buffer 与时钟同样重要}

显示和音频不是“把数据写给设备就结束”。显示控制器按刷新率扫描 framebuffer 或 plane；音频控制器按采样时钟从 ring buffer 取 PCM 数据。它们的核心约束是时间连续性。

\[
audio\ underrun \Leftarrow producer\_rate < consumer\_rate
\]

\[
frame\ tear \Leftarrow scanout(buffer)\ \text{期间修改同一可见区域}
\]

\begin{lstlisting}
display_flip:
  render next buffer
  submit page flip request
  hardware waits for vblank
  scanout switches plane base address
  fence signals old buffer may be reused

audio_playback:
  application fills PCM ring
  driver programs DMA period
  device consumes at sample_rate
  period interrupt wakes refill path
\end{lstlisting}

这里的 completion/fence 不只是性能工具，而是所有权边界。显示旧 buffer 只有在 scanout 不再使用后才能复用；音频 buffer 如果 refill 晚了会 underrun；视频 buffer 如果 fence 没等就写，会造成撕裂或数据竞争。OS 在这些设备上管理的是“内存 + 时间 + 所有权”三件事。

\subsection{TPM、Secure Boot 和 IOMMU：安全硬件如何进入 OS}

安全硬件也要按结构、状态、接口读。TPM 不是“加密芯片”四个字，它有 PCR、密钥、命令队列和访问控制。Secure Boot 不是单个开关，而是从固件到 bootloader 到内核模块的一条验证链。IOMMU 不是只为性能，它限制设备 DMA 范围，防止设备或恶意外设读写任意内存。

\[
measured\ boot:\quad
PCR_{new}=Hash(PCR_{old}\Vert measurement)
\]

\begin{longtable}{p{0.18\linewidth}p{0.34\linewidth}p{0.36\linewidth}}
\toprule
机制 & 硬件状态 & OS 后果 \\
\midrule
TPM PCR & 启动阶段测量累积值 & 密钥封存、远程证明、启动篡改证据 \\
Secure Boot & 固件信任根和签名数据库 & 限制未签名 bootloader/内核模块 \\
IOMMU & 设备 domain 和 IOVA 页表 & DMA 隔离、VFIO 直通、安全边界 \\
CPU 特权/页表 & CPL、PTE、NX、SMEP/SMAP 等 & 用户/内核隔离和执行保护 \\
\bottomrule
\end{longtable}

安全章节如果只讲策略而不讲硬件根，会变成口号。真正的隔离来自多个硬件检查点：CPU 不让 Ring 3 写 CR3，MMU 不让用户访问 supervisor 页，IOMMU 不让设备越界 DMA，TPM 让启动状态可度量。

\subsection{错误报告：硬件失败必须变成可诊断证据}

硬件会失败：ECC corrected error、PCIe AER、IOMMU fault、MCE、NVMe status、USB stall、GPU hang、thermal trip、watchdog reset。高质量 OS 不是假设硬件永远正确，而是把失败变成可定位证据。

\begin{longtable}{p{0.20\linewidth}p{0.34\linewidth}p{0.34\linewidth}}
\toprule
失败源 & 证据 & OS 动作 \\
\midrule
page fault & CR2、error code、RIP、CPL & 修复、发信号、panic 或 oops \\
machine check & bank/status/address/misc & 记录、页下线、进程 kill、系统下线 \\
PCIe AER & requester、status、severity & reset 设备、降级、日志和隔离 \\
IOMMU fault & device/requester、IOVA、权限 & 停止设备、撤销请求、报告驱动 bug \\
NVMe completion error & status code、command id、queue id & 重试、fail request、reset controller \\
thermal/watchdog & trip point、reset reason & 降频、关机、保留 crash 证据 \\
\bottomrule
\end{longtable}

这也给测试方法定了标准。不能只测试正常读写；还要测试超时、reset、IOMMU fault、completion error、热插拔、进程取消、并发关闭、内存压力。硬件路径越异步，错误测试越重要。

\subsection{端到端案例：用户态读文件到 NVMe 完成}

把整节收束成一条路径。用户调用 \code{read(fd, buf, len)}，如果数据不在 page cache，内核可能向 NVMe 提交块 I/O。

\begin{lstlisting}
1. syscall entry saves trap frame and validates fd
2. VFS maps file offset to page cache page
3. filesystem maps logical block to disk LBA
4. block layer builds request and assigns tag
5. DMA API maps page fragments to IOVA
6. NVMe driver fills SQE with opcode, LBA, PRP/SGL
7. dma_wmb establishes descriptor-before-doorbell
8. driver writes queue tail doorbell in BAR
9. controller DMA-reads SQE and DMA-writes data
10. controller writes CQE with command id and status
11. controller sends MSI-X to queue's CPU
12. handler/poll consumes CQE and completes request
13. driver unmaps DMA and block layer unlocks page
14. syscall copies or maps data to user buffer and returns
\end{lstlisting}

每一步对应一个硬件或 OS 对象：

\[
\begin{aligned}
&\text{trap frame}
\rightarrow \text{file/page cache}
\rightarrow \text{bio/request/tag}\\
&\rightarrow \text{IOVA}
\rightarrow \text{SQ/CQ}
\rightarrow \text{MSI-X}
\rightarrow \text{wake up}
\end{aligned}
\]

这条路径能解释很多“为什么”：为什么系统调用要保存寄存器，为什么用户指针不能直接给设备，为什么 page cache 和 CPU cache 不是一回事，为什么 NVMe 有队列和 doorbell，为什么 completion 前不能释放 buffer，为什么中断亲和性会影响吞吐，为什么 \code{fsync} 又比普通 \code{read} 的 completion 多一层持久化语义。

\subsection{本节验收：每个硬件概念都要能落到不变量}

\begin{rubricbox}
\begin{itemize}
\tightlist
\item 给出任意硬件名词，能按“结构、状态、时序、接口、OS 不变量、失败证据”六项拆开。
\item 能解释时钟、复位、电源域为什么决定驱动初始化和 runtime PM 顺序。
\item 能把 \asmcode{add rax, rbx} 追踪到译码、寄存器读、ALU、RFLAGS、提交。
\item 能把 \asmcode{mov rax, qword ptr [rbx + 16]} 追踪到 AGU、TLB、PTE、cache、page fault。
\item 能写出 PTE 权限判定公式，并说明 present、writable、user、NX 的 OS 用途。
\item 能写出 TLB shootdown 的安全释放条件。
\item 能区分 cache coherence、memory ordering、MMIO ordering 和 DMA visibility。
\item 能画出 MSI-X 到 local APIC、IDT、handler、EOI 的中断路径。
\item 能说明 PCIe BAR、bus mastering、DMA mask、IOMMU domain、MSI-X vector 在 probe 顺序中的位置。
\item 能把 NVMe/网卡/virtio 抽象成 submission ring、completion ring、doorbell、DMA buffer 和 request tag。
\item 能说明 USB、I2C、SPI、GPIO 与 PCIe 驱动模型的差别。
\item 能把 TPM、Secure Boot、IOMMU、PTE 这些安全硬件落到具体状态和 OS 后果。
\item 能为 page fault、MCE、AER、IOMMU fault、NVMe error 写出可观测证据。
\end{itemize}
\end{rubricbox}
